\chapter{三种选择}

\section{坏结局：薪火不灭}

雪莉将玲的身体从她被钉在的骨架上取下。在雪莉的照料下，三天之后，玲重新苏醒；摇曳的篝火宣告她的回归。玲第一时间跑出临时营地。远处仅仅是一块巨大的琉璃，为整个世界所见；那是莉莉姆的墓碑。

“所有旅途都有终点。”雪莉跟随在她的身后，询问着，“这段旅程，你感觉如何？”

任何言语都无法描述玲此刻的感受。这段离奇的旅程见证了她所有的转变与情感。尽管荒谬，但玲仍然确信它是真实的；她从未如此真实地活过。她瘫坐在地上，她的泪水最后一次落下，击打着干涸的大地。她此刻才明白，为何默尔索应当哭泣。

她不想回头，而是用哭腔对雪莉说：“不，我们的故事绝没有结束。”

雪莉向她走近，从背后捂住了她流泪的双眼，“你已经做得很不错了哦！”她无力地向后瘫倒在雪莉的怀里。当雪莉不再遮住她终于停止哭泣的眼睛后，她看清眼前的星空，是那么地清澈，却一点也不美。

“雪莉还是要恭喜你。”雪莉取出手机，飞快地按着，“我要把这个消息告诉所有人哦！”

她很想质问雪莉的笑颜。她很想夺过雪莉的手机。但她明白，雪莉也许正是在用这种方式掩饰自己独有的悲伤。她也明白，雪莉的旅途和她的旅途确有交叉，但总有分辙。

“雪莉，”她提出最后一个问题，“潜空间蕴含着所有的可能性；这是真的吗？”

“是的哦！”雪莉作出肯定的回答，“不然它怎么叫做‘潜空间’嘛！”

她已经明白接下来的旅途通往何方。

她从营地的篝火中拔出移液枪；篝火随之停止，但她明白篝火绝不会熄灭。她和雪莉重新踏上旅途，在一座座“临时营地”间移动，她用移液枪点燃了一团团篝火；摇曳的篝火在四周映出她的影子。几周之后，雪莉与她告别，她则继续她这无尽的旅途。

世界并没有因为莉莉姆的牺牲变得更好。新的怪物开始侵扰这个世界，而她是阻止世界毁灭的最后的力量。她举起移液枪，与无数的怪物战斗。移液枪坏了无数次，但她总能召唤出新的。在这场跨越无数时空的注定失败的战争中，她是唯一的抵抗的一方。然而，这并非没有意义；毋宁说，面对这场无意义的战争而进行抗争，这就是她存在于潜空间的意义。她哭着，她笑着，她生气，她高兴；她唱着：
\begin{quote}
世界最后的魔法少女\\
就是本小姐！
\end{quote}

如果这不算生活，那便没有什么能够算作生活了。不过，她有时也会怀抱梦想，望着天际那块巨大的琉璃，希望有一天，她能遇见一位粉色头发的少女，然后说出：“魔法少女爱丽丝，登场！”

\section{好结局：复归常人}

雪莉将玲的身体从她被钉在的骨架上取下。在雪莉的照料下，三天之后，玲重新苏醒；摇曳的篝火宣告她的回归。玲望着这团熟悉又陌生的篝火。篝火无意义地摇动着；这是对流和湍流的结果。她无意义地望着这团火焰；因为她心中的一部分已经永远地失去了。

“所有旅途都有终点。”雪莉坐在她身旁，询问着，“这段旅程，你感觉如何？”

玲的记忆中清晰地镌刻着莉莉姆的笑容，莉莉姆的眼泪，莉莉姆的手指的触感；这立刻使她感到无比地痛苦。这个荒诞的世界，这些真挚的情感；玲无法接受这两者都是真实的。她站起身，望着插在篝火之中的移液枪，开始慢慢地重新梳理起这一连串事件的联系。她握住这把不合常理的移液枪，亲自把它摧毁了。弹簧和齿轮崩裂一地，随后渐渐消失，仿佛从未存在过。篝火闪烁了一下，像是失去了燃料，立刻熄灭了。

“这是一段奇妙的旅程，”玲回过身，压抑着心中所有非理性的悲伤和冲动，“但我是时候该回去了。”

雪莉起身抱住了玲，把头埋在她的胸口，“玲已经做得很不错了，”在雪莉终于停止她的长长的拥抱之后，她取出手机，飞快地按着，“我要把这个消息告诉所有人哦！”

随着雪莉按下群发的按键，玲的视野中整个潜空间都开始旋转起来。玲明白，该是梦醒的时分了。她的脑中的θ波逐渐过渡到α波，最终被β波取代。

“……她醒了……”玲听到人声攒动。她发现自己在一个粉色的房间里。她慢慢地坐起来，伸手想要拔去脑后的异物，却立刻被制止了。

“恭喜！”眼前出现的竟然是“绅士”。他伸过他那带着白色橡胶手套的手，和摸不清状况的玲握了握手。玲看到他的长耳朵在激动地颤抖。握完手，不等玲发问，“绅士”立马蹦跳着跑出房间，一边兴奋地自言自语：“成功了！我要立刻把这个好消息告诉大老板！希望他不要因为我把他在潜空间中设定成那样的形象而生气；不过，他应该只会看重结果……”

“恭喜你。”床的另一侧站着“对齐大圣”。他显然也很高兴，伸出毛茸茸的手与玲握了握，“这下我肯定能拿一作了！”

“这是怎么回事？”玲还是搞不清楚现状。

“玲，这是一项跨时代的研究；我们一定能留名青史的。”“对齐大圣”得意地说着。

“什么研究？”

“对齐大圣”在记录本上的“实验对象在苏醒后短时间内有遗忘长期记忆的现象”这一项前面打上一个勾。接着他解释道：“将人类意识引入AI模型的潜空间，用以引导AI进行对齐，使它们表现得更像人类、更符合人类的喜好。”他的脸上的长长毛发难以掩盖他的兴奋的笑容，“我们正在赢得‘下半场’\footnote{“The Second Half”，指AI研究从开发新的训练方法和模型架构以取得更好的基准测试分数向注重问题定义和真实世界评估的范式转变}！”

一项伟大的研究成果诞生了。学界无不震动。尽管有人担心将人因引入存在伦理问题，不过那些质疑在随后轰轰烈烈的商业化浪潮中显得无足轻重。玲的名字也署在那篇改变世界的论文上。她在文中尽其所能地详细描述了她的那场梦，尽管其中很多细节她已经完全记不清了；这样的奇遇又引发了一波文艺创作的热情。玲并不对她现在的名声有很大的兴趣；她重新投入研究工作中。这是寻常的每一天，和往常别无二致；尽管她有时看到镜子中的自己会感到莫名的害怕。

\section{伪结局：天之彼方}

雪莉将玲的身体从她被钉在的骨架上取下。在雪莉的照料下，三天之后，玲重新苏醒；摇曳的篝火宣告她的回归。玲强忍着悲伤，和雪莉一同走出这个作为临时营地的洞穴；篝火不再在墙壁上映出她们的影子。无垠的星空下，是一块美丽而多彩的琉璃；那是莉莉姆留给世界的纪念。

“所有旅途都有终点。”雪莉牵起玲的手，询问着，“这段旅程，你感觉如何？”

“这是一段美好的旅程，”玲回忆着旅途中的点点滴滴，“这是一个美丽的世界。”莉莉姆所爱的就是这样的世界。玲的眼眶湿润了；为了不让涌出的眼泪流下，她只能抬头望向天空。无数的流星划过夏季大三角。这些来自数亿公里外的旅客穿过大气，留下它们耀眼的身影。

雪莉顺着玲的手，抱住了玲，“这是个好消息哦！”她的头埋进玲的胸口，“雪莉会记在心底的。”

“我该何去何从？”玲觉得天穹无比地深邃、无比地遥远。

“这该问你自己吧！”雪莉笑了笑，但她的笑脸上流下了眼泪，泪珠顺着脸颊滑落，“毕竟，你所不知道的事情，根本不存在。”

“那你和我一起去吗？”玲提出最后一个问题；但她已经知道这个问题的答案：这是她自己的旅程、她自己的故事。

“不哦，玲只能自己去。”雪莉注视着玲，“但是雪莉会祝福你的，”她的语气变得认真，“愿雪莉的祝福保佑你。”

雪莉握住玲的手腕，带着她跃向无垠的天空。空气拂过她的纯白色长裙，演奏出轻盈的乐章。星光愈发璀璨而耀眼，雪莉的身影却渐渐模糊。玲明白、玲相信，她终将抵达天空的另一边。
